% appendix_ladder.tex: the across-borrower term margin adjusted for the penalty and selection, as the check on the within-borrower
% estimate (Code/local/bureau_rho_ladder.R, window_formula.R; exhibits make_ladder_tables.R), the term margin across and
% within borrower by decile (I6f), and the ladder by group. Moved from the main text on 2026-09-16.
\section{The Across-Borrower Term Margin, Adjusted}\label{appendix:ladder}

The bureau's contracts supply both responses the sufficient statistic of Proposition \ref{prop:sensT} needs: a longer contract at the same monthly payment adds payments at the end, and the payment margin moves every payment at once. Their ratio across borrowers, estimated within lender, place and score-by-income cells on 7.6 million auto loans, contains the two channels Proposition \ref{prop:termmargin} names besides the far payments' weight, the balance-scaled penalty and selection into contract length, and Section \ref{sec:termmargin} measures the first. The paper's estimate is the within-borrower one of Section \ref{sec:rate_people}, in which holding the borrower fixed removes selection by design and the penalty channel has decayed by twenty-four months. This appendix is the check: it starts from the across-borrower ratio at twelve months, adjusts for the measured size of each channel one at a time, and reports the rate the window formula returns from the adjusted ratio, with the within-borrower estimate as the reference. The procedure is the one the sufficient-statistics literature follows when a formula derived from a model is evaluated with the best available estimates of its inputs and the identification of each input is stated: \citet{davila_exemptions_2020} derives the optimal bankruptcy exemption as a function of four measurable objects, estimates the sensitivity of credit supply to exemptions from a branch fixed-effects regression on state exemption changes, and reports the sensitivity of the conclusion to the assumed utility parameters.

\subsection{The adjustments}

The outcome is default within twelve months of origination, which removes the exposure term because every loan is observed over the same window. The raw ratio $R_0 = \varepsilon_T/\varepsilon_m$ is $0.01$ for charge-off and $-0.77$ for 90-day delinquency: at face value the far payments have no weight, and the formula returns a rate of one or above. The first adjustment is the penalty channel. A longer contract at the same payment is a larger balance, and default forfeits more of it, which lowers default and biases the term response downward. The deficiency-judgment contrast of Section \ref{sec:termmargin} (Table \ref{tab:term_deficiency}) measures the deficiency part of that channel: where the lender cannot pursue the balance, the term elasticity is higher by $0.36$ (standard error $0.13$) at twelve months. The whole balance channel is $\psi$ times its deficiency part; $\psi = 1$ is the case in which the deficiency judgment is the whole of the balance-scaled cost, and $\psi = 2.5$ allows the credit record's report of the charge-off to add one and a half times as much again, and both are reported because the data do not measure $\psi$. The second adjustment is selection into contract length. Lenders screen with maturity and borrowers sort into it \citep{hertzberg_screening_2018}, so the across-borrower term response differs from the within-borrower one; the within-borrower estimate on multi-loan borrowers is $0.19$ below the same-sample across-borrower estimate at twelve months (Table \ref{tab:within_by_group}), and that difference is applied. The adjusted ratio is mapped to a rate by the window formula of equation \eqref{eq:window} at $K = 12$ with the panel's default timing, $T$ the mean contract length, $h$ the panel's exit hazard and the marginal-utility profile of equation \eqref{eq:lambda_p} at the population's default rate ($\lambda_{60} = 0.34$ for the auto population). The exit hazard here is the panel's gross payoff hazard, $0.22$ per year, whereas the within-borrower estimate of Section \ref{sec:rate_people} uses the hazard net of servicer transfers, $0.19$; with these inputs the formula's value at a rate of zero is $0.43$, and with the within-borrower estimate's it is $0.59$.

Table \ref{tab:bureau_ladder} reports the pooled result. For charge-off with $\psi = 1$, the penalty adjustment alone takes the ratio to $0.81$, above the formula's value at a rate of zero, and the selection adjustment brings it to $0.39$ (standard error $0.31$), which the formula maps to $\rho = 0.04$ with a set from zero to infinity; the set is wide because the two adjustments are measured once, on pooled samples, with standard errors of the order of the ratio itself. The premium's shape is a second-order input here: with a constant premium of $0.2$ at every later date the adjusted ratio maps to $0.15$ and with $0.5$ to $0.18$, inside the same set. With $\psi = 2.5$ the adjusted ratio is above the formula's ceiling and maps to zero with a set bounded at $0.32$. For 90-day delinquency the raw ratio is negative, and the adjustments do not bring it above zero at $\psi = 1$. What the procedure measures with precision is the size of the adjustments a rate taken from contracts across borrowers requires: the penalty channel and selection into contract length together move the term elasticity by more than half a payment elasticity, and the bureau's contracts have no variation that separates them from the far payments' weight. The adjusted across-borrower ratio, $0.04$ with a set that spans the grid, is consistent with the within-borrower estimate of $0.11$ and differs from it in precision: the adjustments cost the across-borrower object the precision that holding the borrower fixed keeps, and the agreement is the check that the penalty channel the within-borrower estimate leaves unsubtracted at twenty-four months is of the size subtracted here at twelve.

\subsection{Selection measured by group}

The selection adjustment above is one number. Table \ref{tab:within_by_group} measures it decile by decile, on the 13.5 million borrowers with two or more auto loans in the ten-percent archive, 65 million loans, within borrower and origination-year cells, by comparing the term elasticity across borrowers within lender-year cells with the same elasticity within borrower, the group being fixed by the borrower's score at first sight so that both loans of a borrower fall in one decile. Across borrowers the term elasticity at twenty-four months climbs from $-0.39$ in the bottom decile to $1.21$ in the top. Within borrower it rises from $0.19$ to $0.43$, and the within-borrower payment elasticity rises with it, from $0.36$ to $0.87$, so the ratio of the two, the object of Proposition \ref{prop:sensT}, is between $0.32$ and $0.53$ in every decile with standard errors of $0.02$ to $0.10$. The across-borrower gradient is selection into contract length, from $+0.58$ in the bottom decile to $-0.78$ in the top, and it has the sign structure of \citet{hertzberg_screening_2018}: positive for subprime borrowers, whom lenders screen with maturity, and negative for prime borrowers, who choose long terms themselves. At twelve months the within-borrower ratio has a gradient, from $-1.58$ in the bottom decile to about $0.3$ in the top three, so the far payments' weight measured within the first year is negative for the borrowers whose balances are charged off most, which is the balance-scaled penalty of Proposition \ref{prop:termmargin}, and the twenty-four-month ratio is where the penalty and the preference channel are both in view. By income quartile, on the 1.35 million multi-loan borrowers with an income at first sight, the twenty-four-month ratio is $0.34$, $0.52$, $0.45$ and $0.27$ from the bottom quartile to the top, with standard errors of $0.05$ to $0.12$, and selection runs from $+0.38$ to $-0.81$: no monotone gradient in income, as in score. The rise of the within-borrower ratio from twelve to twenty-four months, from $-0.70$ to $0.55$ in the pooled sample, each measured to a few hundredths, is a fact about the timing of the far payments' weight that the model of Appendix \ref{appendix:multiperiod} does not reproduce at any rate or penalty (Appendix \ref{appendix:mcliquidity}); it is reported as measured.

\begin{table}[H]
\centering
\caption{The term margin across and within borrower, by score decile}
\label{tab:within_by_group}
\small
\resizebox{\textwidth}{!}{\begin{tabular}{lrrrrrrrrr}
\toprule
 & & \multicolumn{4}{c}{Within twelve months} & \multicolumn{4}{c}{Within twenty-four months} \\
\cmidrule(lr){3-6}\cmidrule(lr){7-10}
Score decile & Loans & Across & Within & Selection & Ratio & Across & Within & Selection & Ratio \\
\midrule
All & 62,518,904 & -0.12 (0.01) & -0.31 (0.01) & -0.19 (0.01) & -0.70 (0.02) & 0.22 (0.00) & 0.30 (0.01) & 0.08 (0.01) & 0.55 (0.01) \\
1 & 6,200,400 & -0.68 (0.01) & -0.45 (0.01) & 0.22 (0.02) & -1.58 (0.07) & -0.39 (0.01) & 0.19 (0.01) & 0.58 (0.01) & 0.53 (0.03) \\
2 & 6,212,641 & -0.54 (0.02) & -0.35 (0.02) & 0.19 (0.02) & -0.95 (0.05) & -0.29 (0.01) & 0.23 (0.01) & 0.52 (0.01) & 0.50 (0.02) \\
3 & 6,223,474 & -0.38 (0.02) & -0.22 (0.02) & 0.16 (0.03) & -0.44 (0.04) & -0.16 (0.01) & 0.27 (0.01) & 0.43 (0.02) & 0.47 (0.02) \\
4 & 6,224,917 & -0.22 (0.02) & -0.13 (0.03) & 0.09 (0.03) & -0.23 (0.04) & 0.02 (0.01) & 0.29 (0.02) & 0.27 (0.02) & 0.44 (0.02) \\
5 & 6,221,828 & -0.08 (0.02) & -0.08 (0.03) & -0.00 (0.04) & -0.13 (0.04) & 0.16 (0.02) & 0.27 (0.02) & 0.11 (0.02) & 0.37 (0.03) \\
6 & 6,239,127 & 0.27 (0.03) & 0.04 (0.04) & -0.23 (0.05) & 0.04 (0.05) & 0.45 (0.02) & 0.27 (0.02) & -0.19 (0.03) & 0.32 (0.03) \\
7 & 6,260,275 & 0.56 (0.04) & 0.05 (0.05) & -0.51 (0.07) & 0.05 (0.06) & 0.78 (0.03) & 0.35 (0.03) & -0.43 (0.04) & 0.39 (0.04) \\
8 & 6,285,765 & 0.96 (0.06) & 0.33 (0.07) & -0.63 (0.09) & 0.31 (0.07) & 1.05 (0.04) & 0.40 (0.04) & -0.65 (0.06) & 0.40 (0.04) \\
9 & 6,315,195 & 1.09 (0.07) & 0.39 (0.09) & -0.70 (0.11) & 0.41 (0.10) & 1.21 (0.05) & 0.41 (0.05) & -0.81 (0.07) & 0.44 (0.06) \\
10 & 6,335,282 & 1.23 (0.09) & 0.32 (0.12) & -0.92 (0.15) & 0.32 (0.12) & 1.21 (0.06) & 0.43 (0.08) & -0.78 (0.10) & 0.49 (0.10) \\
\midrule
\multicolumn{10}{l}{\emph{Income quartile (borrowers with income at first sight)}} \\
1 & 1,541,188 & -0.49 (0.03) & -0.43 (0.03) & 0.06 (0.05) & -1.07 (0.10) & -0.21 (0.02) & 0.17 (0.02) & 0.38 (0.03) & 0.34 (0.05) \\
2 & 1,557,805 & -0.30 (0.04) & -0.27 (0.05) & 0.03 (0.06) & -0.52 (0.10) & -0.02 (0.03) & 0.31 (0.03) & 0.32 (0.04) & 0.52 (0.05) \\
3 & 1,572,950 & 0.45 (0.06) & 0.06 (0.08) & -0.39 (0.10) & 0.08 (0.11) & 0.61 (0.04) & 0.36 (0.05) & -0.25 (0.07) & 0.45 (0.07) \\
4 & 1,578,545 & 0.86 (0.12) & -0.09 (0.16) & -0.96 (0.20) & -0.17 (0.29) & 1.03 (0.08) & 0.22 (0.09) & -0.81 (0.12) & 0.27 (0.12) \\
\bottomrule
\end{tabular}
}
\vspace{0.5em}
\begin{minipage}{0.92\textwidth}
\noindent \small \textit{Notes:} Elasticities of default within twelve and twenty-four months to the contract length at the same payment (with the log payment as the second regressor), on borrowers with two or more auto loans in the ten-percent archive with a credit score at first sight, by the decile of that score. Across: lender-year fixed effects on the same sample; within: borrower and origination-year fixed effects; selection: within minus across, with the standard error of the difference; ratio: the within-borrower term elasticity over the within-borrower payment elasticity, delta-method standard error. Standard errors clustered by borrower.
\end{minipage}
% replication: Code/identification/I6f_within_person_by_group.R, Code/exhibits/make_within_group_table.R
\end{table}

\begin{table}[H]
\centering
\caption{The across-borrower term margin adjusted for the penalty and selection, pooled}
\label{tab:bureau_ladder}
\small
\resizebox{\textwidth}{!}{\begin{tabular}{llrrrrrr}
\toprule
Outcome & $\psi$ & $\varepsilon_m$ & $\varepsilon_T$ & Raw $R$ ($\rho$) & + penalty ($\rho$) & + selection $R$ (s.e.) & $\rho$ [95\% set] \\
\midrule
Charge-off within twelve months & 1.0 & 0.45 (0.02) & 0.01 (0.05) & 0.01 (1.06) & 0.81 (0.00) & 0.39 (0.31) & 0.04 [0.00, $\infty$] \\
Charge-off within twelve months & 2.5 & 0.45 (0.02) & 0.01 (0.05) & 0.01 (1.06) & 2.01 (0.00) & 1.59 (0.73) & 0.00 [0.00, 0.32] \\
90-day delinquency within twelve months & 1.0 & 0.35 (0.02) & -0.27 (0.04) & -0.77 ($\infty$) & 0.25 (0.19) & -0.29 (0.39) & $\infty$ [0.00, $\infty$] \\
90-day delinquency within twelve months & 2.5 & 0.35 (0.02) & -0.27 (0.04) & -0.77 ($\infty$) & 1.78 (0.00) & 1.24 (0.92) & 0.00 [0.00, $\infty$] \\
\midrule
Rule-based experiments, Dobbie--Song pair (for reference) & & & & & & & best fit 0.00; [0.00, 1.20] (5\%), [0.00, 2.84] (1\%) \\
\bottomrule
\end{tabular}
}
\vspace{0.5em}
\begin{minipage}{0.92\textwidth}
\noindent \small \textit{Notes:} Auto loans in the monthly one-percent panel (7.6 million); elasticities of default within twelve months to the payment ($\varepsilon_m$) and to the contract length at the same payment ($\varepsilon_T$), with the cells of Table \ref{tab:main_margins} (lender by origination year, zip3 by origination year, score ventile by income quartile). Raw $R = \varepsilon_T/\varepsilon_m$; the penalty adjustment adds $\psi \times 0.36$ to $\varepsilon_T$ (the deficiency-judgment contrast at twelve months, small-loan specification, ten-percent extract); the selection adjustment adds $-0.19$ (within-borrower minus across-borrower term elasticity at twelve months on the multi-loan borrowers of the ten-percent archive, Table \ref{tab:within_by_group}). The formula maps the adjusted ratio to $\rho$ with the profile of equation \eqref{eq:lambda_p} at the panel's default rate ($\lambda_{60} = 0.34$), the panel's gross payoff hazard $0.22$ (Table \ref{tab:rho_within} uses the hazard net of servicer transfers, $0.19$) and $T = 60$ months; ratios above the formula's value at $\rho = 0$ ($0.43$ at these inputs) map to zero, ratios at or below zero to infinity. Sets: the delta-method standard error of the adjusted ratio, including the deficiency contrast's, taken through the formula. The last row reports the Dobbie--Song pair for reference: its best fit and the sets it accepts at the five and one percent levels (Table \ref{tab:experiments}).
\end{minipage}
% replication: Code/local/bureau_rho_ladder.R, Code/local/window_formula.R, Code/exhibits/make_ladder_tables.R
\end{table}

\subsection{The adjusted ratio by group}\label{appendix:strategies:groupladder}

Table \ref{tab:bureau_rho_groups} applies the same adjustments to every group, with the group's own marginal-utility profile, exit hazard and contract length and the pooled adjustments scaled by the group's charged-off share of the balance. The adjusted ratios are the measured objects, and their ordering is the same at both penalty scalings: the bottom score tercile and the bottom income quartile have negative adjusted ratios, $-0.40$ ($0.29$) and $-0.40$ ($0.32$), which the formula maps to no weight on far payments, while the top two score terciles and the top two income quartiles have adjusted ratios at or above the formula's value at a rate of zero, $0.37$ to $0.43$ at those groups' inputs, which it maps to full weight. Applied group by group the procedure returns boundary values of the mapping, ratios measured to about $\pm 0.3$; the ordering across groups is the exposure of Section \ref{sec:heterogeneity} seen through the term margin, since the payment elasticity in the denominator rises with the score and with income, and it is not a gradient in the rate, which the within-borrower design of Section \ref{sec:rate_people} measures directly and finds flat.

\begin{table}[H]
\centering
\caption{The adjusted across-borrower ratio by group}
\label{tab:bureau_rho_groups}
\small
\resizebox{\textwidth}{!}{\begin{tabular}{llrrrrrrrr}
\toprule
Dimension & Group & Loans & $\lambda_T$ & $\varepsilon_m$ & $\varepsilon_T$ & Raw $R$ & Corrected $R$ (s.e.) & $\rho$ [set], $\psi = 1$ & $\rho$ [set], $\psi = 2.5$ \\
\midrule
All &  & 7,355,247 & 0.34 & 0.45 & 0.01 & 0.01 & 0.39 (0.31) & 0.04 [0.00, $\infty$] & 0.00 [0.00, 0.32] \\
Score tercile & 1 & 2,472,131 & 0.38 & 0.48 & -0.37 & -0.77 & -0.40 (0.29) & $\infty$ [0.36, $\infty$] & 0.00 [0.00, $\infty$] \\
Score tercile & 2 & 2,455,327 & 0.31 & 0.79 & 0.46 & 0.58 & 0.74 (0.18) & 0.00 [0.00, 0.00] & 0.00 [0.00, 0.00] \\
Score tercile & 3 & 2,427,789 & 0.26 & 0.47 & 1.31 & 2.81 & 2.95 (0.65) & 0.00 [0.00, 0.00] & 0.00 [0.00, 0.00] \\
Income quartile & 1 & 1,846,938 & 0.38 & 0.45 & -0.37 & -0.83 & -0.40 (0.32) & $\infty$ [0.27, $\infty$] & 0.00 [0.00, $\infty$] \\
Income quartile & 2 & 1,851,364 & 0.34 & 0.57 & 0.13 & 0.22 & 0.47 (0.23) & 0.00 [0.00, 1.06] & 0.00 [0.00, 0.13] \\
Income quartile & 3 & 1,845,440 & 0.30 & 0.82 & 0.58 & 0.71 & 0.83 (0.17) & 0.00 [0.00, 0.00] & 0.00 [0.00, 0.00] \\
Income quartile & 4 & 1,811,505 & 0.27 & 0.74 & 0.63 & 0.84 & 0.96 (0.25) & 0.00 [0.00, 0.00] & 0.00 [0.00, 0.00] \\
Region & MW & 1,598,359 & 0.34 & 0.50 & 0.04 & 0.07 & 0.44 (0.30) & 0.00 [0.00, $\infty$] & 0.00 [0.00, 0.26] \\
Region & NE & 1,160,318 & 0.33 & 0.67 & 0.28 & 0.42 & 0.63 (0.23) & 0.00 [0.00, 0.31] & 0.00 [0.00, 0.00] \\
Region & S & 2,962,049 & 0.35 & 0.39 & -0.10 & -0.25 & 0.18 (0.35) & 0.27 [0.00, $\infty$] & 0.00 [0.00, $\infty$] \\
Region & W & 1,581,037 & 0.34 & 0.49 & -0.06 & -0.13 & 0.23 (0.29) & 0.20 [0.00, $\infty$] & 0.00 [0.00, 1.11] \\
Cohort & 2005-09 & 1,995,599 & 0.31 & 0.54 & -0.35 & -0.65 & -0.45 (0.27) & $\infty$ [0.57, $\infty$] & 0.06 [0.00, $\infty$] \\
Cohort & 2010-12 & 1,226,517 & 0.33 & 0.38 & -0.21 & -0.55 & -0.27 (0.36) & $\infty$ [0.00, $\infty$] & 0.00 [0.00, $\infty$] \\
Cohort & 2013-15 & 1,077,245 & 0.35 & 0.41 & 0.01 & 0.02 & 0.41 (0.39) & 0.03 [0.00, $\infty$] & 0.00 [0.00, 0.67] \\
Cohort & 2016-18 & 1,108,633 & 0.35 & 0.44 & 0.13 & 0.29 & 0.68 (0.45) & 0.00 [0.00, $\infty$] & 0.00 [0.00, 0.14] \\
Cohort & 2019-21 & 1,045,442 & 0.35 & 0.38 & 0.04 & 0.12 & 0.67 (0.47) & 0.00 [0.00, $\infty$] & 0.00 [0.00, 0.17] \\
Cohort & 2022+ & 901,811 & 0.36 & 0.39 & 0.06 & 0.16 & 0.81 (0.48) & 0.00 [0.00, $\infty$] & 0.00 [0.00, 0.02] \\
\bottomrule
\end{tabular}
}
\vspace{0.5em}
\begin{minipage}{0.92\textwidth}
\noindent \small \textit{Notes:} Charge-off within twelve months; the adjustments of Table \ref{tab:bureau_ladder} with the group's own marginal-utility profile ($\lambda_T$ is its weight on the last payment; Appendix \ref{appendix:mucorrection}), exit hazard and mean contract length, and the pooled adjustments scaled by the group's charged-off share of the balance. Cohorts are origination years.
\end{minipage}
% replication: Code/local/bureau_rho_ladder.R, Code/exhibits/make_ladder_tables.R
\end{table}
